% Drafted against audit/outline.md §4, with one addition: a subsection
% specifying the age-weighted quantile. The outline predates the discovery
% that Goldman's disclosed weighting had never been implemented, so its
% bullets do not include it, and the estimator has to be defined formally
% before §7 can report the sweep that turns on it.
%
% Figures from data/processed/book_fixed_notionals.csv and config/params.yaml,
% hash-recorded in audit/evidence_freeze.csv at commit 1e6bdd74.

\section{Method}\label{sec:method}

\subsection{Instruments, and what full revaluation means here}

Equity, currency and commodity positions are linear in their factor: a one-day
fractional return on the factor is the fractional profit and loss on the
position. Interest-rate positions are not. They are repriced by discounting
their cashflows again at the shocked yield, which is what Goldman's phrase
``full valuation'' describes, rather than multiplied by a modified duration.

The two rate sub-sleeves keep their source's native convention. US positions
are par coupon bonds priced off the Treasury par yield curve with semiannual
coupons; euro positions are zero-coupon bonds priced off the ECB spot curve
with annual compounding. The alternative would be to convert one curve into the
other's convention, which needs a bootstrap the published data does not
support at these tenors. The cost of keeping both is that a dollar of euro zero
and a dollar of US par bond are not interchangeable: a zero has longer duration
than a par bond of the same maturity, so the sleeves differ per unit of
notional. The sizing described below targets risk rather than notional, which
absorbs the difference, but the instrument sets are not identical and
Section~\ref{sec:limitations} records that.

Writing $P(y, m)$ for the price of the relevant instrument at yield $y$ and
maturity $m$, the scenario return for a rate position under a historical yield
change $\Delta y_{s}$ observed on day $s$ is
\begin{equation}
r_{s} = \frac{P(y_{t} + \Delta y_{s},\, m)}{P(y_{t},\, m)} - 1 ,
\label{eq:revaluation}
\end{equation}
where $y_{t}$ is the yield at the close of the valuation date. The level enters
as well as the change, which is why the engine carries both: the same 25 basis
point move is worth a different amount of money at a 0.5\% yield than at a 5\%
one. Equation~\eqref{eq:revaluation} is also where convexity enters, and
Section~\ref{sec:assumptions} reports what it is worth by running the same
scenarios through the first-order approximation
$r_{s} \approx -D_{\mathrm{mod}}(y_{t}, m)\,\Delta y_{s}$ instead.

\subsection{Sizing within a sleeve, and sizing the sleeve against the disclosure}

Sizing happens in two stages, and only the second touches the disclosures.

Within a sleeve, instruments are weighted so that each contributes equal
standalone VaR over the calibration window. This is a neutrality choice rather
than an empirical claim: nothing is disclosed about the composition of
Goldman's rate book, so loading the sleeve onto any one tenor would be an
assumption with no support. Equal risk contribution is the weighting that
requires no such assumption.

The sleeve as a whole is then scaled so that its average daily VaR across the
calibration year equals the figure Goldman disclosed for that category. The
scaling is available in closed form, which is worth stating because it is the
design of the experiment rather than a computational convenience. Profit and
loss on every position here is proportional to its notional, so for a sleeve
$\cat$ with notionals $\mathbf{n}$ scaled by $k > 0$,
\begin{equation}
\VaR_{\cat}(k\mathbf{n}) = k \, \VaR_{\cat}(\mathbf{n}) ,
\label{eq:homogeneity}
\end{equation}
and the same proportionality carries through the average over the year. One
division per sleeve therefore lands exactly on the disclosed target.
Section~\ref{sec:validation} confirms that it does, to machine precision.

Equation~\eqref{eq:homogeneity} is positive homogeneity, which is one of the
four axioms \citet{artzner1999coherent} use to define a coherent risk measure
and one of the three VaR actually satisfies. It is worth naming because the
axiom VaR fails --- subadditivity --- is the subject of this paper's central
quantity, and the calibration leans on the one property that holds while the
experiment measures the one that does not.

The consequence is that the summed category VaR matches the disclosures by
construction and carries no information whatever. That is deliberate. Fixing
the four category figures leaves total VaR and hence the diversification effect
entirely free, because the total depends on how the four sleeves move together
and nothing in the calibration constrains that. Comparing the free quantity
against its disclosed counterpart is the experiment; the matched quantities are
the control.

Table~\ref{tab:book} gives the resulting book. Its dollar level is a
normalisation and should not be read as an estimate of anything: only the
ratios between sleeves carry meaning, and those come from the disclosures.

% Generated by build/make_tables.py from book_fixed_notionals.csv. Do not edit.

\begin{table}[tbp]
\centering
\small
\caption{The fixed proxy book after calibration to the FY2019 disclosed risk split}
\label{tab:book}
\begin{tabular}{@{}lllrr@{}}
\toprule
Factor & Category & Instrument & Maturity (y) & Notional (\$mm) \\
\midrule
us\_2y & Interest rates & Par coupon bond & 2 & 8,403 \\
us\_5y & Interest rates & Par coupon bond & 5 & 2,467 \\
us\_10y & Interest rates & Par coupon bond & 10 & 1,302 \\
us\_30y & Interest rates & Par coupon bond & 30 & 553 \\
eur\_2y & Interest rates & Zero-coupon bond & 2 & 18,162 \\
eur\_5y & Interest rates & Zero-coupon bond & 5 & 3,941 \\
eur\_10y & Interest rates & Zero-coupon bond & 10 & 1,427 \\
\addlinespace
equity\_us & Equity prices & Linear in factor & --- & 604 \\
equity\_europe & Equity prices & Linear in factor & --- & 478 \\
equity\_japan & Equity prices & Linear in factor & --- & 444 \\
equity\_em & Equity prices & Linear in factor & --- & 442 \\
\addlinespace
fx\_eurusd & Currency rates & Linear in factor & --- & 788 \\
fx\_usdjpy & Currency rates & Linear in factor & --- & 734 \\
fx\_gbpusd & Currency rates & Linear in factor & --- & 701 \\
\addlinespace
commodity\_gold & Commodity prices & Linear in factor & --- & 443 \\
commodity\_oil & Commodity prices & Linear in factor & --- & 165 \\
commodity\_base\_metals & Commodity prices & Linear in factor & --- & 341 \\
\bottomrule
\end{tabular}
\notetext{Notionals are the output of the sizing described in Section~\ref{sec:method}, not an estimate of anything Goldman holds. Sleeve totals are \$36,255mm of rates, \$1,968mm of equity, \$2,223mm of currency and \$949mm of commodity exposure, \$41,395mm in all. The rate sleeve dominates on notional and not on risk: a two-year note carries a fraction of the daily risk of a dollar of crude oil fund, which is the whole reason sizing targets risk rather than money. The dollar level is a normalisation and carries no information; only the ratios between sleeves do.}
\end{table}


\subsection{The engine and its timing convention}

For a valuation date $t$, the scenario set is the $\window$ most recent daily
factor moves up to and including $t$, applied to the book as it stands at the
market levels of $t$. Each scenario produces one profit-and-loss outcome per
sleeve and one for the book as a whole, and VaR is the loss at the
$\cover$ quantile of that distribution, reported as a positive number. Category
VaR uses only that category's positions. Total VaR moves every factor together,
historical day by historical day, which is what preserves the co-movement the
paper is testing. The diversification effect is
\begin{equation}
\divef_{t} = \VaR_{t}^{\text{total}} - \sum_{\cat \in \cats} \VaR_{t}^{\cat} ,
\qquad
\divpct_{t} = \frac{-\divef_{t}}{\sum_{\cat \in \cats} \VaR_{t}^{\cat}} ,
\label{eq:diversification}
\end{equation}
following Goldman's definition exactly, with $\divpct_{t}$ the share reported
throughout this paper.

The settings are the disclosed ones where the filings supply them:
$\window = 1{,}260$ trading days, which is the five years Goldman states it
samples from, and $\cover = 0.05$ for the 95\% risk-management measure. The
quantile is read off the finite scenario set by linear interpolation; the
discrete order statistic is the conservative alternative and
Section~\ref{sec:assumptions} reports the difference rather than assuming it
away. Appendix~\ref{sec:appendix_b} gives the reading conventions in full,
together with the argument behind the closed-form calibration and the
look-ahead probe.

The timing convention is stated explicitly because getting it wrong produces a
model that looks better than it is. VaR dated $t$ uses only the book, the
levels and the scenarios available at the close of $t$, and it is compared
against the profit and loss realised on $t+1$. No quantity dated $t+1$ enters
the VaR dated $t$. This is Goldman's own convention: the filings state that
daily net revenues are compared with VaR calculated as of the end of the prior
business day. Section~\ref{sec:validation} describes the test that
enforces it.

\subsection{Scenario age weighting}

Goldman states that its historical data is weighted so that the importance of
older observations decays. Implementing that replaces the equally weighted
quantile with a weighted one, following the hybrid approach of
\citet{boudoukh1998hybrid}. A scenario observed $k$ trading days before the
valuation date enters with weight
\begin{equation}
w_{k} = \frac{\lambda^{k}}{\sum_{j=0}^{\window-1} \lambda^{j}} ,
\qquad \lambda \in (0, 1] ,
\label{eq:ageweights}
\end{equation}
and VaR is read at the point where cumulative weight, rather than cumulative
count, reaches $\cover$. At $\lambda = 1$ every weight is $1/\window$ and the
estimator returns to the equally weighted percentile.

Two properties of the implementation matter for interpreting
Section~\ref{sec:assumptions}. The first is that the collapse at $\lambda = 1$
is exact, in both the interpolated and the discrete reading, which
Section~\ref{sec:validation} reports as a test: any difference between a
weighted and an unweighted run is therefore the weighting and not a change of
estimator. The second is that $\lambda$ is not disclosed. The filings describe
the scheme but not its rate, so the baseline keeps $\lambda = 1$, leaving no
undisclosed parameter inside it, and the rate is swept instead, at values
chosen for their half-lives rather than for how close they land to the
disclosures.

\subsection{Two books, and what each is for}

Two books are run, and conflating them would be a mistake.

The \emph{recalibrated} book is resized every year to that year's own disclosed
split. Its exposure mix tracks the disclosures exactly, which removes the
exposure mix from the comparison and leaves only the cross-category
correlations in the public data. It is the right book for asking whether the
diversification effect comes out where Goldman reports it. It is not a
forecast, and could not be: sizing the 2022 book uses a figure published in
February 2023.

The \emph{fixed} book is sized once, on 2019, and then left alone. Nothing
about it responds to anything disclosed afterwards, so the path its VaR takes
is driven purely by what happened to volatility and correlation. It is the
right book for the regime comparison, and the only defensible one for
backtesting, because its notionals were set before any of the data it is scored
on. 2019 is the base year because it is calm and precedes everything the
backtest scores, which makes 2020 onward out of sample with respect to sizing.

\subsection{The volatility-weighted variant}

One variant is run alongside the disclosed specification, because it addresses
a failure the disclosed specification has.
\citet{hullwhite1998} rescale each historical return by the ratio of today's
conditional volatility to the volatility in force when the return was observed,
so the scenario set is restated in current volatility units before the
percentile is taken. Volatility is estimated by exponentially weighted variance
with the \citet{riskmetrics1996} daily decay factor of 0.94, and the estimator
is causal by construction: the forecast stamped $t$ uses only moves up to $t$.
A GARCH(1,1) alternative is refitted each January on the history available at
that point.

Realised profit and loss always comes from the unscaled moves. What the book
actually made on $t+1$ is not a function of which model forecast it, and
scaling the outcome as well as the scenarios would make the backtest circular.

Age weighting and volatility weighting are easy to confuse and do different
things. Age weighting changes how much a scenario counts according to when it
happened. Volatility weighting changes the size of the scenario itself
according to the volatility regime it came from. Both turn out to repair the
engine's coverage, and neither repairs the diversification result, for a reason
Section~\ref{sec:assumptions} develops: both act on the scale of the scenario
set, and the diversification effect is a function of its correlation structure.
